\documentclass[12pt,a4paper]{article}
\usepackage{mathwizards-formulario}
\title{\textbf{Integrales Dobles — Ejercicios Resueltos}}
\author{}
\date{}

% ── Comandos propios de este documento ─────────
\newcommand{\resultado}[1]{%
  \textcolor{resultcolor}{\textbf{Resultado: } $\displaystyle #1$}%
}

\begin{document}
\maketitle

\begin{center}
\Large\bfseries Parte I — Integrales Dobles Indefinidas
\end{center}

%---------------------------------------------------------------------------
% EJERCICIO 1 (Indefinida)
%---------------------------------------------------------------------------
\section*{Ejercicio 1}
\[
  \iint (3x^2 + 2x^3)^2 \, dy \, dx
\]

\textbf{Paso 1 — Expandir el cuadrado.}
\[
  (3x^2 + 2x^3)^2 = 9x^4 + 12x^5 + 4x^6
\]

\textbf{Paso 2 — Integrar respecto a $y$ (la variable interna).}

El integrando solo depende de $x$, así que $y$ actúa como factor lineal:
\[
  \int (9x^4 + 12x^5 + 4x^6)\,dy = (9x^4 + 12x^5 + 4x^6)\,y + C_1(x)
\]

\textbf{Paso 3 — Integrar respecto a $x$.}
\[
  \int \left[(9x^4 + 12x^5 + 4x^6)\,y\right] dx
  = y \left(\frac{9x^5}{5} + \frac{12x^6}{6} + \frac{4x^7}{7}\right) + C_2(y)\,x + C_3
\]
\[
  = y \left(\frac{9x^5}{5} + 2x^6 + \frac{4x^7}{7}\right) + C_2(y)\,x + C_3
\]

\medskip
\resultado{\dfrac{9x^5 y}{5} + 2x^6 y + \dfrac{4x^7 y}{7} + C_2(y)\,x + C_3}

\bigskip
%---------------------------------------------------------------------------
% EJERCICIO 2 (Indefinida)
%---------------------------------------------------------------------------
\section*{Ejercicio 2}
\[
  \iint (\sqrt{2x} - \sqrt{5y} + 3) \, dx \, dy
\]

\textbf{Paso 1 — Integrar respecto a $x$ (variable interna).}

Reescribimos: $\sqrt{2x} = \sqrt{2}\,x^{1/2}$.
\[
  \int \left(\sqrt{2}\,x^{1/2} - \sqrt{5y} + 3\right) dx
  = \sqrt{2}\cdot\frac{x^{3/2}}{3/2} - \sqrt{5y}\,x + 3x + C_1(y)
\]
\[
  = \frac{2\sqrt{2}}{3}\,x^{3/2} - x\sqrt{5y} + 3x + C_1(y)
\]

\textbf{Paso 2 — Integrar respecto a $y$.}

Integramos término a término. Para $\sqrt{5y} = \sqrt{5}\,y^{1/2}$:
\[
  \int \left(\frac{2\sqrt{2}}{3}\,x^{3/2} - x\sqrt{5}\,y^{1/2} + 3x\right) dy
\]
\[
  = \frac{2\sqrt{2}}{3}\,x^{3/2}\,y - x\sqrt{5}\cdot\frac{y^{3/2}}{3/2} + 3x\,y + C_1(y)\text{ (integrada)} + C_2
\]
\[
  = \frac{2\sqrt{2}}{3}\,x^{3/2}\,y - \frac{2x\sqrt{5}}{3}\,y^{3/2} + 3xy + C
\]

\medskip
\resultado{\dfrac{2\sqrt{2}}{3}\,x^{3/2}\,y - \dfrac{2\sqrt{5}}{3}\,x\,y^{3/2} + 3xy + C}

\bigskip
%---------------------------------------------------------------------------
% EJERCICIO 3 (Indefinida)
%---------------------------------------------------------------------------
\section*{Ejercicio 3}
\[
  \iint (3x^{-1} + 2y^{-2}) \, dy \, dx
\]

\textbf{Paso 1 — Integrar respecto a $y$ (variable interna).}
\[
  \int (3x^{-1} + 2y^{-2})\,dy = 3x^{-1}\,y + 2\cdot\frac{y^{-1}}{-1} + C_1(x)
  = \frac{3y}{x} - \frac{2}{y} + C_1(x)
\]

\textbf{Paso 2 — Integrar respecto a $x$.}
\[
  \int \left(\frac{3y}{x} - \frac{2}{y}\right) dx
  = 3y\ln|x| - \frac{2x}{y} + C
\]

\medskip
\resultado{3y\ln|x| - \dfrac{2x}{y} + C}

\bigskip
%---------------------------------------------------------------------------
% EJERCICIO 4 (Indefinida)
%---------------------------------------------------------------------------
\section*{Ejercicio 4}
\[
  \iint \left(\frac{2y}{\sqrt{y}} + \frac{2x}{\sqrt{x}}\right) dx \, dy
\]

\textbf{Paso 1 — Simplificar el integrando.}
\[
  \frac{2y}{\sqrt{y}} = 2y^{1/2}, \qquad
  \frac{2x}{\sqrt{x}} = 2x^{1/2}
\]
El integrando queda $2\sqrt{y} + 2\sqrt{x}$.

\textbf{Paso 2 — Integrar respecto a $x$ (variable interna).}
\[
  \int (2\sqrt{y} + 2x^{1/2})\,dx
  = 2\sqrt{y}\,x + 2\cdot\frac{x^{3/2}}{3/2} + C_1(y)
  = 2x\sqrt{y} + \frac{4}{3}\,x^{3/2} + C_1(y)
\]

\textbf{Paso 3 — Integrar respecto a $y$.}
\[
  \int \left(2x\,y^{1/2} + \frac{4}{3}\,x^{3/2}\right) dy
  = 2x\cdot\frac{y^{3/2}}{3/2} + \frac{4}{3}\,x^{3/2}\,y + C
\]
\[
  = \frac{4x}{3}\,y^{3/2} + \frac{4}{3}\,x^{3/2}\,y + C
\]

\medskip
\resultado{\dfrac{4x}{3}\,y^{3/2} + \dfrac{4}{3}\,x^{3/2}\,y + C}

\bigskip
%---------------------------------------------------------------------------
% EJERCICIO 5 (Indefinida)
%---------------------------------------------------------------------------
\section*{Ejercicio 5}
\[
  \iint \left(\frac{y\sqrt{2}}{x} + \frac{x\sqrt{2}}{y}\right)^2 dy \, dx
\]

\textbf{Paso 1 — Expandir el cuadrado.}

Sea $a = \dfrac{y\sqrt{2}}{x}$ y $b = \dfrac{x\sqrt{2}}{y}$. Entonces $(a+b)^2 = a^2 + 2ab + b^2$:
\[
  a^2 = \frac{2y^2}{x^2}, \qquad
  b^2 = \frac{2x^2}{y^2}, \qquad
  2ab = 2\cdot\frac{y\sqrt{2}}{x}\cdot\frac{x\sqrt{2}}{y} = 2\cdot 2 = 4
\]
\[
  \text{Integrando} = \frac{2y^2}{x^2} + 4 + \frac{2x^2}{y^2}
\]

\textbf{Paso 2 — Integrar respecto a $y$ (variable interna).}
\[
  \int \left(\frac{2y^2}{x^2} + 4 + \frac{2x^2}{y^2}\right) dy
  = \frac{2}{x^2}\cdot\frac{y^3}{3} + 4y + 2x^2\cdot\frac{y^{-1}}{-1} + C_1(x)
\]
\[
  = \frac{2y^3}{3x^2} + 4y - \frac{2x^2}{y} + C_1(x)
\]

\textbf{Paso 3 — Integrar respecto a $x$.}
\[
  \int \left(\frac{2y^3}{3x^2} + 4y - \frac{2x^2}{y}\right) dx
  = \frac{2y^3}{3}\cdot\frac{x^{-1}}{-1} + 4xy - \frac{2}{y}\cdot\frac{x^3}{3} + C
\]
\[
  = -\frac{2y^3}{3x} + 4xy - \frac{2x^3}{3y} + C
\]

\medskip
\resultado{-\dfrac{2y^3}{3x} + 4xy - \dfrac{2x^3}{3y} + C}

\newpage
\begin{center}
\Large\bfseries Parte II — Integrales Dobles Definidas
\end{center}

%---------------------------------------------------------------------------
% EJERCICIO 6 (Definida 1)
%---------------------------------------------------------------------------
\section*{Ejercicio 6}
\[
  \int_{0}^{2} \int_{1}^{2} \left(\frac{7}{2x} - \frac{5}{6y}\right) dy \, dx
\]

\textbf{Nota:} $\sqrt{4} = 2$, por lo que el límite superior de la integral interna es $2$.

\textbf{Paso 1 — Integral interna (respecto a $y$, de $1$ a $2$).}
\[
  \int_{1}^{2} \left(\frac{7}{2x} - \frac{5}{6y}\right) dy
  = \left[\frac{7y}{2x} - \frac{5}{6}\ln|y|\right]_{1}^{2}
\]
\[
  = \left(\frac{7\cdot 2}{2x} - \frac{5}{6}\ln 2\right) - \left(\frac{7}{2x} - 0\right)
  = \frac{14}{2x} - \frac{7}{2x} - \frac{5}{6}\ln 2
\]
\[
  = \frac{7}{2x} - \frac{5\ln 2}{6}
\]

\textbf{Paso 2 — Integral externa (respecto a $x$, de $0$ a $2$).}

\textbf{Atención:} $\dfrac{7}{2x}$ tiene una singularidad en $x = 0$. Evaluamos como impropia:
\[
  \int_{0}^{2} \left(\frac{7}{2x} - \frac{5\ln 2}{6}\right) dx
  = \lim_{\varepsilon \to 0^+} \left[\frac{7}{2}\ln|x| - \frac{5\ln 2}{6}\,x\right]_{\varepsilon}^{2}
\]
\[
  = \lim_{\varepsilon \to 0^+} \left[\left(\frac{7}{2}\ln 2 - \frac{5\ln 2}{3}\right) - \left(\frac{7}{2}\ln \varepsilon - \frac{5\ln 2}{6}\,\varepsilon\right)\right]
\]
Como $\ln\varepsilon \to -\infty$ cuando $\varepsilon\to 0^+$, la integral \textbf{diverge}.

\medskip
\textcolor{red}{\textbf{La integral diverge}} (singularidad en $x = 0$).

\bigskip
%---------------------------------------------------------------------------
% EJERCICIO 7 (Definida 2)
%---------------------------------------------------------------------------
\section*{Ejercicio 7}
\[
  \int_{-4}^{-3} \int_{-1}^{1} (7x^2 - 12x^3 y^6) \, dx \, dy
\]

\textbf{Paso 1 — Integral interna (respecto a $x$, de $-1$ a $1$).}
\[
  \int_{-1}^{1} (7x^2 - 12x^3 y^6)\,dx
  = \left[\frac{7x^3}{3} - \frac{12x^4 y^6}{4}\right]_{-1}^{1}
  = \left[\frac{7x^3}{3} - 3x^4 y^6\right]_{-1}^{1}
\]

Evaluamos en $x=1$:
\[
  \frac{7}{3} - 3y^6
\]
Evaluamos en $x=-1$:
\[
  -\frac{7}{3} - 3y^6
\]
Restamos:
\[
  \left(\frac{7}{3} - 3y^6\right) - \left(-\frac{7}{3} - 3y^6\right) = \frac{14}{3}
\]

\textbf{Paso 2 — Integral externa (respecto a $y$, de $-4$ a $-3$).}
\[
  \int_{-4}^{-3} \frac{14}{3}\,dy = \frac{14}{3}\left[y\right]_{-4}^{-3}
  = \frac{14}{3}(-3 - (-4)) = \frac{14}{3}\cdot 1 = \frac{14}{3}
\]

\medskip
\resultado{\dfrac{14}{3}}

\bigskip
%---------------------------------------------------------------------------
% EJERCICIO 8 (Definida 3)
%---------------------------------------------------------------------------
\section*{Ejercicio 8}
\[
  \int_{-2}^{2} \int_{-1}^{0} \left(\frac{2y}{\sqrt{x}} + \frac{6x}{\sqrt{9}}\right) dx \, dy
\]

\textbf{Paso 1 — Verificar dominio.}

El término $\dfrac{2y}{\sqrt{x}}$ requiere $x > 0$. Sin embargo, la integral interna va de $x = -1$ a $x = 0$, donde $\sqrt{x}$ no está definida en los reales.

\medskip
\textcolor{red}{\textbf{La integral no existe}} (integrando no definido en $\mathbb{R}$ para $x < 0$).

\bigskip
%---------------------------------------------------------------------------
% EJERCICIO 9 (Definida 4)
%---------------------------------------------------------------------------
\section*{Ejercicio 9}
\[
  \int_{-3}^{-2} \int_{-2}^{0} \left(\frac{3 - x - 6y}{\sqrt{9x}}\right) dy \, dx
\]

\textbf{Paso 1 — Verificar dominio.}

$\sqrt{9x}$ requiere $x \geq 0$, pero la integral externa va de $x = -3$ a $x = -2$, donde $9x < 0$.

\medskip
\textcolor{red}{\textbf{La integral no existe}} ($\sqrt{9x}$ no está definida para $x < 0$ en $\mathbb{R}$).

\bigskip
%---------------------------------------------------------------------------
% EJERCICIO 10 (Definida 5)
%---------------------------------------------------------------------------
\section*{Ejercicio 10}
\[
  \int_{1/2}^{1} \int_{0}^{2} (2x^{-1} + 3y^{-2}) \, dx \, dy
\]

\textbf{Paso 1 — Verificar singularidades.}

$2x^{-1} = \dfrac{2}{x}$ tiene una singularidad en $x=0$, que es el límite inferior de la integral interna. Evaluamos como impropia.

\textbf{Paso 2 — Integral interna (respecto a $x$, de $0$ a $2$), impropia en $x=0$.}
\[
  \lim_{\varepsilon\to 0^+}\int_{\varepsilon}^{2} \left(\frac{2}{x} + 3y^{-2}\right) dx
  = \lim_{\varepsilon\to 0^+}\left[2\ln|x| + 3y^{-2}\,x\right]_{\varepsilon}^{2}
\]
\[
  = \lim_{\varepsilon\to 0^+}\left[\left(2\ln 2 + 6y^{-2}\right) - \left(2\ln\varepsilon + 3y^{-2}\varepsilon\right)\right]
\]

Como $\ln\varepsilon \to -\infty$ cuando $\varepsilon \to 0^+$, la integral interna \textbf{diverge}.

\medskip
\textcolor{red}{\textbf{La integral diverge}} (singularidad en $x = 0$).

\end{document}
